\chapter{Vrange problem}

\begin{motto}
Faget tok det som skulle vore objektet for ein teori,\\ og gjorde det til ei orsaking for å sleppe å lage teorien.
\end{motto}

Då metoderørsla kollapsa, sat designfaget att med eit problem det aldri har kome seg unna sidan: korleis grunngje at det \emph{ikkje} kan ha eit rasjonelt fundament, utan å innrømme at det dermed knapt er eit fag. Løysinga kom utanfrå, frå byplanlegginga, i form av eit omgrep som skulle vise seg å bli det mest verdfulle retoriske verktøyet faget nokon gong har lånt: \emph{vrange problem}. Med det kunne faget omkode si djupaste veikskap til sin eigentlege natur. Dette kapitlet sporar omgrepet frå opphavet hjå Horst Rittel til kanoniseringa hjå Richard Buchanan, og viser korleis ein presis og rett observasjon vart vridd til eit alibi.

\section{Rittel og Webber: ein presis observasjon}

I 1973 publiserte Horst Rittel og Melvin Webber ein artikkel om dilemma i ein generell planleggingsteori \autocite{rittel1973}. Poenget deira var skarpt og korrekt. Problema byplanleggjaren står overfor, skreiv dei, liknar ikkje problema vitskapsmannen eller ingeniøren løyser. Dei er \emph{vrange} («wicked», i motsetnad til «tame»): dei har inga endeleg formulering, ingen stoppregel som seier når dei er løyste, inga rett-eller-gale-løysing men berre betre eller verre, og kvart slikt problem er i praksis unikt og kan ikkje testast ved prøving og feiling, fordi kvart forsøk endrar situasjonen irreversibelt. Rittel og Webber retta dette mot den teknokratiske optimismen i etterkrigsplanlegginga — mot førestillinga, beslekta med Herbert Simons rasjonalisme, om at samfunnsproblem kunne løysast med same systematiske metode som tekniske problem.

Observasjonen var verdfull, og ho gjeld utan tvil for mykje av det designarar gjer. Å forme ein offentleg teneste, ein by, eit grensesnitt som tusenvis skal leve med, er ikkje å løyse eit veldefinert problem med ein kjend prosedyre. Så langt er Rittel og Webber sine allierte, ikkje motstandarar, av denne boka. Spørsmålet er ikkje om designproblem er vrange. Spørsmålet er kva ein gjer med den innsikta.

\section{Buchanan: frå observasjon til identitet}

Det avgjerande steget tok Richard Buchanan i 1992, i ein av dei mest siterte artiklane i heile designteorien \autocite{buchanan1992wicked}. Buchanan henta vrange problem inn i design og gjorde dei til sjølve hjørnesteinen i ein teori om designtenking. Design, hevda han, er ikkje knytt til eit bestemt fagområde slik medisinen er knytt til kroppen; design er ein allmenn evne til å gje form, og emnet for design er nett det ubestemte, det vrange. Designaren si oppgåve er å gje førebels form til problem som per natur manglar endeleg form.

Det er ein elegant tanke, og han forklarte mykje av kvifor designfaget kjendest annleis enn nabofaga. Men sjå kva det gjorde med fundamentspørsmålet. Frå Rittel sin observasjon — \emph{nokre problem er vrange} — til Buchanan sin posisjon — \emph{designens vesen er å handtere det vrange} — er det eit stort steg, og steget har ein pris. For om det vrange er designens eigentlege element, så følgjer det at design ikkje \emph{kan} ha eit fundament av det slaget dei tamme faga har: ingen kumulative resultat, ingen overførbar metode, ingen falsifiserbare påstandar. Veikskapen er ikkje lenger ein mangel; han er identiteten. Spør kvifor designforskinga ikkje akkumulerer slik biologien gjer, og svaret er no innebygd: fordi problema er vrange, og vrange problem produserer ikkje kumulativ kunnskap. Kritikken har avvæpna seg sjølv før han er framført.

\section{Forvekslinga: objekt mot orsaking}

Her ligg den djupaste feilen i resonnementet, og ho er verdt å seie heilt presist. At eit problem er vrangt, tyder ikkje at det ikkje kan finnast ein teori \emph{om kvifor} det er vrangt, og \emph{korleis} ein navigerer det. Tvert om: at formgjeving må balansere mange samtidige, motstridande krav over eit høgdimensjonalt moglegheitsrom, er ikkje ein grunn til at teori er umogleg. Det er innhaldet i den teorien som trengst.

Samanlikninga som avslører dette er biologien. Ein levande organisme løyser i kvart augneblink eit problem like vrangt som noko ein designar møter: den må tilfredsstille utalige motstridande krav samstundes — energi, struktur, forsvar, reproduksjon — utan nokon stoppregel, utan ei rett løysing, i ein situasjon som endrar seg irreversibelt. Evolusjonsbiologien stod overfor nett denne vrangskapen, og svara ikkje med å erklære seg sjølv uvitskapleg. Han bygde eit presist apparat for nett denne situasjonen: tilpassingslandskapet, der mange samtidige seleksjonstrykk summerer seg til ein topografi som forma navigerer \autocite{wright1932}. Vrangskapen vart ikkje ein grunn til å gje opp teorien; ho vart spesifikasjonen teorien måtte oppfylle.

Designfaget hadde det same valet og tok det motsette. Det tok vrangskapen som prov på at det ikkje kunne ha eit fundament, i staden for som krava til fundamentet. Og det er ikkje slik at alternativet var utenkjeleg på Buchanans tid: Herbert Simon hadde alt formulert «ill-structured problems» som noko ein kunne studere strukturen i, ikkje berre bøye seg for \autocite{simon1973illstructured}. Skilnaden mellom Simon og Buchanan er skilnaden mellom å spørje «kva er strukturen i dårleg strukturerte problem?» og å svare «dei har ingen, og det er vår fridom».

\section{Alibiet i bruk}

Når omgrepet fyrst var på plass, vart det faget si vitskapsteoretiske immuncelle. Kvar gong nokon spør designforskinga om etterprøvbar, kumulativ kunnskap, kan ho minne spørjaren om at problema er vrange, og at kravet difor er malplassert — ein kategorifeil, ein som forvekslar design med naturvitskap. Det er eit effektivt svar, fordi det inneheld ein sann kjerne. Men effekten er at faget har gjort seg sjølv uangripeleg på den eine måten som betyr noko: det kan ikkje lenger ta feil. Eit rammeverk som spesifiserer kva som ville tilbakevise det, kan vere vitskap; eit som har bygd inn at ingenting kan tilbakevise det, er noko anna. Christopher Alexander, som hadde prøvd å gje design ein eksplisitt metode og så forkasta sitt eige prosjekt, såg faren frå den andre kanten \autocite{alexander1971state}: han forlét metoden fordi han hadde blitt eit sjølvgåande spel med diagram, lausrive frå om bygningane vart gode. Mellom Alexanders desillusjon og Buchanans omfamning av det vrange ligg heile faget si manglande evne til å halde på eit fundament — anten det forkastar metoden fordi han er tom, eller hyllar formløysa fordi ho er fri.

Resultatet ser ein i litteraturen. Etter eit halvt århundre med «wicked problems» som ramme finst det ingen veksande kropp av etablerte designresultat som studentane lærer som grunnlag, slik biologistudenten lærer seleksjon. Det finst ein veksande bunke av einskildstudiar, kvar med sin eigen vinkel, ingen avgjord sann eller usann. Eit felt der ingen påstand nokon gong vert avgjort, nærmar seg ikkje sanninga. Det samlar meiningar med fotnotar. Og det gjer det med god samvit, fordi det har lært seg å kalle fråvêret av kumulativ kunnskap for respekt for problemets natur.

\vignett

Den same logikken — at design er for djupt, for situert, for vrangt til å la seg formalisere — strukturerer òg faget si forklaring på kvar kunnskapen då eigentleg bur. Svaret er: i den reflekterande utøvaren, i atelieret, i blikket. Det er det neste kapitlet.

\vidarelesing{\fullcite{rittel1973}; \fullcite{buchanan1992wicked}; \fullcite{simon1973illstructured}; \fullcite{alexander1971state}; \fullcite{cross2011dt}.}
